% ------------------------------------------------------------------------
% file `main-bookexercise-1-solution-body.tex'
%
%     solution of type `bookexercise' with id `1'
%
% generated by the `exerciseanswer' environment of the
%   `xsim' package v0.21 (2022/02/12)
% from source `main' on 2026/10/07 on line 769
% ------------------------------------------------------------------------
1. D。线性表要求元素存在位置上的逻辑关系，A错误；线性表是一种逻辑结构，通常可以用数组或链表实现，B错误；线性表要求同一个表中的数据类型相同，C错误；实际需求中我们
要处理int型的数据，就会使用一个int型的线性表，要处理Role类型的数据，就会使用Role类型的线性表，D正确。\\

2. C。A、B选项都需要移动若干元素，时间复杂度为O(n)；D项需要从表头开始，将每一个元素的值与x对比直至找到，时间复杂度为O(n)。C项，由于每个元素
的地址可以由其位序直接计算得来，因此时间复杂度为O(1)。\\

3. C。\\

4. B。由于链表结构中保存有表首元素指针，因此在表头插入元素时间复杂度一定是O(1)。对于C选项，要在单链表的表尾插入元素，需要先遍历链表，拿到尾节点的指针，再
操作链条，时间复杂度为O(n)。我们将在后面看到，为了解决单链表无法高效在尾部插入元素的问题，可以通过单链表的多种变态解决，如：记录尾指针的
\uwave{循环单链表}、\uwave{循环双链表}等。\\

5. 事实上，这种存储方式下，还需要一个字段capacity标记当前已动态分配的空间总量：
\begin{cppcode}
    struct SeqList {
        ElemType* data;     // 数组起始地址
        int current_size;  // 当前表中的元素数量
        int capacity;      // 当前已分配的数组容量
    };
\end{cppcode}
例如：capacity = 100，current\_size = 80表示分配了100个位置的内存空间，目前使用了80个，还剩20个。操作方面，显然，初始化动态分配版本的顺序表需要
为data分配一个初始空间，初始分配多大你可以自己定义。除初始化外，这里我们仅指出一个重点操作：插入。

回忆静态存储的顺序表，若表满，只能拒绝插入（当然你可以规定淘汰算法，接受新值淘汰旧值，这取决于你，但静态存储已经无法存储更多了）。而动态存储版本，
当已分配空间满了的时候，我们可以申请一块更大的内存空间，将已有表中数据整体拷贝过去，并接受新元素。C++中的std::vector就采用了类似机制实现。\\

6. 拟采用带头节点的实现。代码如下：

1）要构建\(a_1\) → \(a_2\) → ... → \(a_{10}\)，不难想出，只需要按顺序将每一个元素依次链接到链表尾部就可以，代码如下：
\begin{cppcode}
    LinkList f1(int a[], int n = 10) {
        // 创建头结点
        LinkList head = (int*)malloc(sizeof(ListNode));
        head->next_node = NULL;
        // 动态维护tail，让其始终指向当前链表最后一个节点
        LinkList tail = head;
        for (int i = 0; i < n; i++) {
            LinkList node = (int*)malloc(sizeof(ListNode));   // 为新数据分配一个节点
            node->data = a[i];
            node->next_node = NULL;
            // 新节点插入在最后一个节点之后（尾插）
            tail->next_node = node;
            tail = node;    // 动态维护tail指针，让tail始终指向最后一个节点
        }
        return head;
    }
\end{cppcode}

    2）要构建\(a_{10}\) → \(a_9\) → ... → \(a_1\)，实际也很简单，只需要按顺序将每一个元素依次插入链表头部即可，代码如下：
    \begin{cppcode}
        LinkList f2(int a[], int n = 10) {
            // 创建头结点
            LinkList head = new ListNode;
            head->next_node = NULL;
            for (int i = 0; i < n; i++) {
                LinkList node = new ListNode;
                node->data = a[i];
                // 新节点插入在头结点之后（头插）
                node->next_node = head->next_node;
                head->next_node = node;
            }
            return head;
        }
    \end{cppcode}
    事实上，函数f1就是许多经典教材中\term{“尾插法”}的实现，函数f2就是\term{“头插法”}的实现。\\


7. 思路：遍历一遍链表，“检查”每一个节点的值，若等于e，将其删除，并重构该节点附近链条。代码如下。
\begin{cppcode}
    void DeleteValue(LinkList L, ElemType e) {
        ListNode* pre = L;          // 当前“检查”的节点的前驱节点
        ListNode* cur = L->next_node; // 当前“检查”的节点
        while (cur != NULL) {
            if (cur->data == e) {
                ListNode* temp = cur;
                // 重构链条
                pre->next_node = cur->next_node;
                // 推移cur指针，准备检查下一个节点
                cur = cur->next_node;
                // 释放删除节点
                free(temp);
            }
            else {
                // 当前节点不是目标节点
                // 前驱和当前节点同时后移, 继续检查下一个节点
                pre = cur;
                cur = cur->next_node;
            }
        }
    }
\end{cppcode}

8. 相信大家很自然地就能想到的一种做法是：先遍历一遍链表，计算出表长n，再从头开始查找\(\frac{n+1}{2}\)位置处的节点将其返回。代码如下：
\begin{cppcode}
    NODE* FindMiddle1(NODE *L) {
        int n = 0;
        // 第一次遍历，统计节点数量
        NODE *p = L->next;  // 带头结点，L->next是第一个数据节点
        while(p != NULL) {
            n++;
            p = p->next;
        }
        // 空表情况
        if(n == 0)
            return NULL;
        // 第二次遍历，寻找第(n+1)/2个节点
        int pos = (n + 1) / 2;
        p = L->next;    // p指针重置为第一个数据节点，从这里开始找
        for(int i = 1; i < pos; i++) {
            p = p->next;
        }
        return p;
    }
\end{cppcode}

事实上，还有一种更高效的算法，就是所谓的经典的\term{“双指针算法”}。我们设置两个指针，slow指针每次走一步，fast指针每次走两步，这样，slow“走过的里程”永远是fast的
一半，所以，当fast走到表尾时，slow就处在中间。需要注意的就是边界判断及奇偶的处理。代码如下：
\begin{cppcode}
    NODE* FindMiddle2(NODE *L) {
        NODE *slow = L->next;
        NODE *fast = L->next;
        if(slow == NULL)
            return NULL;
        while(fast->next != NULL && fast->next->next != NULL) {
            slow = slow->next;
            fast = fast->next->next;
        }
        return slow;
    }
\end{cppcode}

9. 本题思路很简单。对于带头节点的单链表，我们依次将（本来就在头节点后的\(a_1\)、）\(a_2\)、\(a_3\)、...、\(a_n\)插入到第一个数据节点之前（即头节点之后）就行。
我们动态维护指针p，让其指向当前“待插”节点。代码如下。
\begin{cppcode}
    NODE* Reverse(NODE* L) {
        NODE* p = L->next;  // L是头节点，L->next是第一个数据节点
        L->next = NULL;     // 插入过程中，让链表的“已定形部分”的尾节点指向NULL
        while(p) {
            // next为下一个待插节点
            NODE* next = p->next;
            // 插入当前节点
            p->next = L->next;
            L->next = p;
            // 动态维护p，使其指向下一个待插节点
            p = next;
        }
        return L;
    }
\end{cppcode}
这个过程其实也可以理解为按顺序将当前链表节点用头插法{\color{red}\textbf{（见习题1 第6小题）}}“重新插入该链表”。

对于不带头结点的单链表同样是将每一个节点依次插入到第一个数据节点之前：
\begin{cppcode}
    NODE* Reverse(NODE* L) {
        NODE *new_head = NULL;
        NODE *p = L;
        while(p != NULL) {
            NODE *next = p->next;
            p->next = new_head;
            new_head = p;
            p = next;
        }
        return new_head;
    }
\end{cppcode}

10. 经过观察分析不难看出，题目要我们将“倒数第i个元素”插入到“第i个元素”后面，换句话说，如果我们将链表从正中间“砍成两半”，题目要求的正是我们将链表的“后半部分”
从后往前依次插入“前半部分”中。大家可能会想，后半部分如果有反向链条就好了对吧？就可以依次一个一个插入了。那没有，怎么办呢？
相信聪明的你已经想到了，在第上一个习题中，我们介绍了链表的倒置算法。我们将后半部分链表倒置一下不就可以了？没错！接下来我们就给出本题的完整思路：

我们先用双指针算法（见习题1 第8小题）找到链表的中点，将链表砍成两半。然后将后半部分链表逆置，再依次插入前半部分的对应位置中。

若元素个数为奇数个，那“中点”应该分给第一部分还是第二部分呢？其实经过简单的思考不难发现：其实给哪部分都行，我们是可以通过代码控制的。

完整的代码实现如下：
\begin{cppcode}

void Rearrange(NODE *head) {
    if(head == NULL || head->next == NULL)
        return;

    // 双指针寻找链表中点，slow最终指向第一部分的最后一个节点
    NODE *slow = head;
    NODE *fast = head;
    while(fast->next != NULL &&
          fast->next->next != NULL) {
        slow = slow->next;
        fast = fast->next->next;
    }
    /* 将链表“砍成两半”
     * 第一部分:
     * head->next ... slow
     * 第二部分:
     * second-> ...
     */
    NODE *second = slow->next;  // second即第二部分的“第一个节点”。注意：第二部分事实上是一个“不带头结点”的单链表，其第一个节点就是second指向的数据节点。
    slow->next = NULL;
    /*
     * 第二步：
     * 逆置第二部分链表
     */
    NODE *prev = NULL;
    NODE *cur = second;
    while(cur != NULL) {
        NODE *next = cur->next;
        cur->next = prev;
        prev = cur;
        cur = next;
    }
    second = prev;
    /*
     * 第三步：
     * 交替合并两个链表
     */
    NODE *first = head->next;
    NODE *p = first;    // 动态维护p，让其始终指向第一部分的插入位置
    NODE *q = second;   // 动态维护q，让其始终指向第二部分的待插节点
    while(p != NULL && q != NULL) {
        NODE *pNext = p->next;
        NODE *qNext = q->next;
        // 插入q节点
        p->next = q;
        q->next = pNext;
        // 推移p和q两个指针，准备插入下一个节点
        p = pNext;
        q = qNext;
    }
}
\end{cppcode}
寻找中点、逆置第二部分链表和交替合并的过程时间复杂度均为O(n)，因此整个算法时间复杂度为O(n)。





